\section{Experiments}
\label{sec:experiments}

We validate the formula-guided onset search of \S\ref{sec:algorithm} and demonstrate it on a
worked scenario. These are \emph{numerical} experiments: they evaluate the algorithm against
the analytic queueing model of Ramani and Tantawi~\cite{ramani2026queueing} --- the same oracle
used throughout the paper --- not a live server. Their role is exploration and validation. The
substantive comparative study --- the end-to-end benefit of capping \emph{concurrency} (serving
at the onset $M^\ast$) against other admission or rate-limiting schemes on a real deployment ---
is future work (\S\ref{sec:exp-threats}).

\subsection{Setup}
\label{sec:exp-setup}
One \emph{oracle} call is one analytic SLO-feasible-rate evaluation $f(M)$ (\S\ref{sec:analysis}).
We evaluate on the six development scenarios spanning the four regimes and on a $30$-scenario
benchmark drawn by Latin-hypercube sampling (seed $42$) over the service- and workload-parameter
ranges; $25$ are feasible and $5$ are infeasible at every $M$ (all-zero $f$ curves), on which a
correct strategy returns $m_{\min}$. For every scenario we precompute a \emph{truth cache}: the
oracle evaluated at all $M\in[1,256]$, giving the ground-truth curve, the strict argmax
$M_{\mathrm{truth}}$, and the peak. The $\varepsilon$-onset $M^\ast$ is the smallest $M$ with
$f(M)\ge(1-\varepsilon)\max_M f(M)$, $\varepsilon=0.02$.\footnote{The Analysis and Algorithm
sections tabulate $M^\ast$ as $M_{\mathrm{truth}}$ (the strict argmax) where the two coincide to
floating point; the genuine $\varepsilon$-onset can sit a few steps below (baseline: $66$ vs
$69$) and far below on float-wiggly benchmark plateaus --- precisely why we report both gaps.}
Because the cache is the oracle tabulated at every $M$ and the strategies are deterministic, we
\emph{replay} the three committed strategies offline against the caches with a counted,
cache-backed oracle (no server); call counts and chosen $M$ reproduce a live run exactly. As in
the harness, one confirmatory $f(\hat M)$ call follows each strategy, so reported call counts
include that $+1$.

\subsection{Onset accuracy and cost}
\label{sec:exp-accuracy}
Table~\ref{tab:eval-comparison} compares the formula-guided search against the two baselines of
\S\ref{sec:algo-baselines} on the feasible benchmark. We report $\mathrm{gap}_{\mathrm{onset}}
=|\hat M - M^\ast|$ against the $\varepsilon$-onset --- the quantity the algorithm optimises ---
and, for continuity with the discovery campaign, $\mathrm{gap}_{\mathrm{argmax}}=|\hat M -
M_{\mathrm{truth}}|$ against the strict argmax. Formula-guided is the only strategy that is at
once \emph{cheap} ($\le 8$ oracle calls), \emph{SLO-feasible} (worst $\mathrm{gap}_f=0.0138<
\varepsilon$), and \emph{near-onset} (worst $\mathrm{gap}_{\mathrm{onset}}=38$, mean $7.6$). The
naive-max heuristic spends a single call but over-provisions massively (mean
$\mathrm{gap}_{\mathrm{onset}}=180$) and even violates the latency tolerance on an interior-peak
scenario ($\mathrm{gap}_f=0.0286$); parameter-blind ternary spends $30$ calls to converge to the
plateau interior, also violating the tolerance and landing far from the onset. Since both
baselines breach $\mathrm{gap}_f\le\varepsilon$, formula-guided is the only SLO-feasible
operating point of the three, and it reaches it at near-onset concurrency in a handful of
calls. The six development scenarios corroborate (worst
$\mathrm{gap}_f=0.0101$, $\le 8$ calls).

\begin{table}[t]
  \centering
  \caption{Onset-search accuracy and cost over the 25 feasible benchmark scenarios (worst / mean). $\mathrm{gap}_{\mathrm{onset}}=|\hat M - M^\ast|$ is measured against the $\varepsilon$-onset (the objective); $\mathrm{gap}_{\mathrm{argmax}}$ against the strict throughput argmax (campaign convention; its floor is structural, RP-10). $\mathrm{gap}_f$ is the relative throughput shortfall; the SLO tolerance is $\varepsilon=0.02$. Only Formula-guided is simultaneously cheap, SLO-feasible, and near-onset.}
  \label{tab:eval-comparison}
  \begin{tabular}{lrrrr}
    \toprule
    Strategy & calls & $\mathrm{gap}_{\mathrm{onset}}$ (worst/mean) & $\mathrm{gap}_f$ (worst) & $\mathrm{gap}_{\mathrm{argmax}}$ (worst) \\
    \midrule
    Formula-guided & 8 & 38/7.6 & 0.0138~\checkmark & 72 \\
    Parameter-blind ternary & 30 & 155/69.6 & 0.0286~$\times$ & 145 \\
    Naive-max & 1 & 251/180.2 & 0.0286~$\times$ & 251 \\
    \bottomrule
  \end{tabular}
\end{table}


The residual $\mathrm{gap}_{\mathrm{argmax}}$ (worst $72$) is dominated by the structural
argmax-vs-onset divergence (\S\ref{sec:search}, RP-10): on float-wiggly plateaus the strict
argmax sits far above the onset (one scenario: argmax $118$ vs onset $44$), so most of that gap
is not search error but the deliberate choice to stop at the onset rather than over-provision to
the argmax.

\subsection{A measured search trace}
\label{sec:exp-trace}
Figure~\ref{fig:onset-search} shows the procedure as actually executed on the baseline scenario:
$f^\ast$ anchored at $m_{\max}=256$, the threshold at $(1-\varepsilon/2)f^\ast$, and the
downward monotone-predicate probes $\{148,94,67,53,60,64\}$ descending from the bracket
$[40,256]$ to $\hat M=67$, one step above the onset $M^\ast=66$ --- eight oracle calls including
the anchor and the confirmatory evaluation.

\subsection{Scope and threats to validity}
\label{sec:exp-threats}
The oracle is the analytic model, not a live server: results inherit its Markovian-arrival and
three-parameter service assumptions, and a real vLLM deployment adds stochastic token lengths,
KV-cache eviction, and chunked prefill that can break the $nc=1$ regime or the M/M/1 wait model.
Calibrating the primitives $\mathrm{itl}(B)$, $\mathrm{ttft\_prefill}(B)$ against measured vLLM
is the headline next step. The exact closed-form $M_{\mathrm{ITL}}$ holds under $nc=1$ (all dev
scenarios; most of the benchmark); under chunk-count jumps the warm start can err by a step
(\S\ref{sec:mitl}). Finally, the evaluation establishes that the algorithm \emph{finds} the
near-minimal SLO-feasible concurrency; quantifying the operational \emph{benefit} of serving
there, versus alternative limiting schemes, is the comparative study left to future work.
